% Nouveau EGA V (cote 25): archive cote 25, site page 133 (PDF page 134).
% typescript. Draft compared: none.
% AI-transcribed by Claude Opus 5.5 (claude-opus-5-5, adaptive thinking, effort high), Message Batches API, batch ega25b (msgbatch_01A7n5YuXubUPNK1zsepi8RG), 2026-11-13; not yet checked by eye.
\scanpage{25/133}
\mspage{36}

un sous-k-espace vectoriel\inserted{E}de l'espace des \deleted{\inserted{pseudo}} fonctions \deleted{rationelles} \inserted{\unsure{méromorphes}} \deleted{\inserted[typed]{strictes}} \xout{xxxxxxx}
sur X, tel que pour tout $\phi \hand{\in} E-(0)$, $\phi$ soit \inserted{\unsure{régulier}} \deleted{inversible dans l'\xout{xxxxxxx}}
\deleted{anneau des fonctions rationelles} \deleted{\inserted[typed]{strictes}} \xout{xxxxxx}, i.e. div($\phi$) est défini \deleted{comme}
\deleted{diviseur de Cartier,}\unsure{\hand{,}} et que $\hand{\triangle}$ soit l'ensemble des div($\phi$) pour $\phi \hand{\in} E-(0)$.
c) Soient E, E' deux sous-\inserted[typed]{k-}espaces vectoriels de l'espace des
fonctions rationelles \inserted[typed]{strictes} \xout{xxxxxxxx} sur X, \xout{xxxxx} satisfaisant la condition
envisagée dans b), alors les ensembles de diviseurs $\hand{\triangle}$, $\hand{\triangle}'$ qu'ils
définissent respectivement sont égaux sss il existe une \deleted{\inserted{\unsure{pseudo-}}} fonction
\deleted{rationelle stricte}e \deleted{inversible} \inserted{\unsure{régulière}} $\phi$ sur X, telle que $E' = \phi E$ . Si $E\hand{\neq}(0)$\inserted{(i.e. $\triangle \neq \emptyset$)}, une
telle $\phi$ est déterminée modulo multiplication par un élément c de $k^{\hand{*}}$.

La démonstration est un exercice facile, enutilisant \hand{\S}.1. , et je
me dispense d'en écrire la démonstration, sauf réclamation de ta part. Il
\margin{||} me sembled'ailleurs que \hand{\S}.4. viendrait avantageusement dès avant \hand{\S}.3.,
étanttechniquement nettement plus trivial. Noter d'ailleurs que lorsque
X est géom intègre, la condition sur E énoncée dans b) devient \uline{vide}.
La restriction mise au début de b) tient au fait qu'il se pourrait autre-%
ment que la condition énoncée pour b) ne soit plus vérifiée après passage
à la clôture algébrique de k (il est facile d'en donner des exemples\xout{xx},
en toute caractéristique, et même lorsque k est séparablement clos en
c\xout{\unsure{ar}} p $>$ 0). Pour bien faire, on devrait énoncer b) sans condition supplé-%
mentaire sur k ou X, mais en énoncant la condition sur E en passant à la
clôture algébrique de k (et signaler que si X est géom intègre, \unsure{oette|cette}
condition devient vide). \hand{-- Par abus de langage, un}
\hand{\uline{ens.} $\triangle$ de diviseurs de la forme \unsure{Ens}(D) sera}
\hand{encore parfois appelé système linéaire de diviseurs}
\hand{sur X ;}

\mspage{133}
